% ATENÇÃO - veja com o seu orientador se você vai ter este capítulo e se este vai ter nome!
\chapter{Trabalhos Relacionados}
\label{cap:trabalhos:relacionados}

Esta seção apresenta a revisão de literatura sobre o tema deste trabalho, com foco em \gls{llms}, fine-tunning e explicabilidade.

\section{Labirintos como Domínio Estrutural para Raciocínio Espacial}

A investigação do raciocínio espacial em modelos de linguagem tem utilizado labirintos como um domínio estruturado e simbólico, capaz de revelar limitações e padrões de generalização dos modelos. O benchmark introduzido por \cite{zhou2024navigatingshortcutmazecomprehensive}, avalia a suscetibilidade dos \gls{llms} a se apoiarem em atalhos (\textit{shortcuts}) superficiais, ou correlações onde não há relação de causa-efeito, que afetam a robustez dos modelos em tarefas de classificação de texto.

Em outra abordagem, \cite{hsu2025whatmakesmazelooklike} propuseram o framework \textit{\gls{dsg}}, que utiliza \gls{llms} para extrair e estruturar representações simbólicas (\textit{schemas}) do conceito de labirinto, promovendo um raciocínio mais alinhado à cognição humana ao combinar essas abstrações com modelos multimodais. Essas contribuições reforçam a relevância dos labirintos textuais como ambiente de análise para o comportamento de \gls{llms} em tarefas que exigem planejamento, abstração e robustez interpretativa.

\section{Raciocínio Espacial em LLMs}

O raciocínio espacial em \gls{llms} tem sido explorado por diferentes abordagens, buscando adaptar esses modelos para tarefas que exigem compreensão e manipulação de ambientes estruturados.

O AlphaMaze \cite{dao2025alphamaze} propõe um processo em duas etapas: inicialmente, o modelo é ajustado via \gls{sft} para aprender comandos de movimentação em labirintos textuais; em seguida, utiliza-se o \gls{grpo} para aprimorar o raciocínio e a autocorreção. O MazeBench é empregado para avaliar o desempenho dos modelos em diferentes níveis de dificuldade, evidenciando ganhos após o \gls{grpo}.

Em paralelo, \cite{wang2024pictureworththousandwords} questionam a vantagem de modelos multimodais em tarefas espaciais, mostrando que \gls{llms} baseados apenas em texto podem superar alternativas multimodais, o que reforça a relevância de investigar estratégias específicas para \gls{llms} textuais. Em contextos tridimensionais, \cite{zhang2025doespointcloudboost} identificam limitações dos \gls{llms}, enquanto \cite{cheng2024spatialrgpt} propõem integrar módulos de representação de regiões e informações de profundidade para melhorar a compreensão espacial.

Outras abordagens, como o \textit{\gls{cos}} \textit{Prompting} \cite{hu2024chainofsymbol}, convertem descrições em linguagem natural para representações simbólicas intermediárias, otimizando o planejamento espacial e reduzindo o uso de tokens. Já o \textit{\gls{vot}} \cite{wu2024visualizationofthought} estimula a geração de representações visuais em \gls{ascii} \textit{art}, inspirando-se na cognição humana para melhorar o acompanhamento de estados e o planejamento de ações. Embora o \gls{vot} apresente avanços em determinados cenários, sua eficácia depende das habilidades emergentes dos modelos e da qualidade dos \textit{prompts}, sendo limitado em tarefas mais complexas ou com modelos menos robustos.

Ainda, mesmo com avanços recentes, variantes como os \textit{\gls{lrms}}, que incorporam mecanismos explícitos de raciocínio, como cadeias de pensamento, \textit{\gls{cot}} \cite{wei2023chainofthought}, e auto-reflexão, ainda enfrentam desafios em tarefas de alta complexidade, apresentando limitações na generalização e fenômenos como raciocínio redundante ou ineficiente \cite{shojaee2025illusionthinking}.

Por fim, no contexto deste trabalho, o raciocínio espacial é central, pois o objetivo é investigar como \gls{llms} podem ser ajustados e analisados para resolver labirintos representados em \gls{ascii}. As abordagens discutidas na literatura fornecem subsídios para a escolha de técnicas de fine-tuning e avaliação, além de motivar o uso de ferramentas de explicabilidade para compreender como as representações internas dos modelos evoluem ao longo do treinamento.

\section{Fine-tuning de LLMs para Tarefas Específicas}

O fine-tuning de \gls{llms} para tarefas específicas é um tema amplamente investigado, com diferentes estratégias sendo propostas para aprimorar a adaptação dos modelos. Por exemplo, \cite{jiang2024supervised} demonstra que o uso de \gls{sft} pode melhorar tanto as representações internas quanto a capacidade de generalização, mesmo em modelos visuais, evidenciando a importância do ajuste direcionado para tarefas de raciocínio espacial.

Além disso, \cite{hsieh2023distillingstepbystep} propõem a utilização de \textit{rationales} gerados por \gls{llms} como supervisão adicional em um cenário multitarefa, combinando rótulos e explicações para enriquecer o sinal de treinamento. Essa abordagem, baseada em \gls{cot} \textit{prompting} \cite{wei2023chainofthought}, permite que modelos menores alcancem desempenho superior a \gls{llms} maiores, mesmo com menos dados, ao tornar o processo de ajuste mais eficiente e interpretável. No contexto deste trabalho, a integração de racionalizações se mostra relevante para aprimorar a eficiência e a interpretabilidade dos modelos.

\section{Explicabilidade em LLMs}

A explicabilidade em \gls{llms} é um campo em expansão, com abordagens que buscam tornar os modelos mais transparentes e compreensíveis. A biblioteca \gls{llm-mri} \cite{costa2024llmmri} destaca-se por permitir a análise e visualização das ativações neuronais, facilitando a interpretação de como os modelos se adaptam a tarefas específicas.

Complementarmente, \cite{wu2025usablexai10strategies} propõem o conceito de \textit{Usable} \gls{xai}, que enfatiza a aplicação prática das explicações para diagnóstico e aprimoramento dos modelos. Entre as estratégias discutidas estão métodos de atribuição, análise das ativações internas e explicações baseadas em exemplos, todas relevantes para entender a evolução das representações internas durante o \textit{fine-tuning}. A integração dessas abordagens contribui para identificar limitações, corrigir vieses e aprimorar o desempenho dos modelos em tarefas que exigem alto grau de interpretabilidade.
